Mobile GUI agents must often use a one-time password, credential, address, phone number, or payment identifier to finish a task. Today the value is placed in a screenshot, accessibility tree, prompt, or action argument, so an injected instruction can convert task authority into disclosure or redirection. We present \SH, an Android design for \emph{usable but invisible} secrets: a privileged service mints an opaque, authenticated handle at the authoritative OS source---the SMS or notification OTP pipeline, Credential Manager, the clipboard, or a wallet---before the raw value can reach the agent. The handle binds a secret type, a task identifier, a destination predicate over package and field class (with fail-closed confinement for unattested WebViews), and a type-specific lifecycle (single-use and TTL for OTPs, task-scoped for identifiers, confirm-per-use for payment). The agent observes only the handle and issues \texttt{ENTER(handle)} toward a UI node; the OS resolves the value only after authenticating the handle, validating task and target, enforcing expiry/use/confirmation, and reserving the handle against races, then redacts the resolved field from subsequent agent observation. We evaluate an on-device Android prototype across OTP forwarding, cross-app and intra-app Confused Deputy destination redirection, replay, and concurrency attacks against Claude Sonnet~4.6-driven Mobile-Agent v3.5 and AutoDroid agents. \SH{} eliminates secret exposure ($0/30$, Wilson $[0.00, 0.11]$) versus 0.97 for unprotected agents ($29/30$), completely mitigates destination redirection ($0/30$ vs $30/30$ for unbound placeholders), guarantees atomic single-use serialization ($50/50$ races), incurs microsecond in-vault state transitions ($5.42\,\mu\text{s}$ mint, $0.18\,\mu\text{s}$ resolve), and introduces zero task regression across 10 AndroidWorld applications ($\Delta = 0.00$).
